% LSPV-1 REPORT::CIRCEQUIV1::2026-10-12T11:10:41Z
% !TEX program = xelatex
\documentclass[11pt]{article}
\usepackage[margin=0.82in]{geometry}
\usepackage{amsmath,amssymb,amsthm,array,booktabs,longtable,xurl,hyperref}
\hypersetup{colorlinks=true,urlcolor=blue,linkcolor=blue}
\newtheorem{theorem}{Theorem}
\newtheorem{lemma}{Lemma}
\newcommand{\poly}{\mathrm{poly}}
\newcommand{\NP}{\mathsf{NP}}
\newcommand{\PP}{\mathsf{P}}
\begin{document}

\section*{LSPV-1 universal-core lower bound audit: the unrestricted-circuit endpoint is not a weaker bridge}

\noindent\textbf{Run ID:} \texttt{CIRCEQUIV1}.\quad
\textbf{Timestamp:} 2026-10-12T11:10:41Z.\quad
\textbf{Overall status:} \textbf{OPEN}.  P versus NP remains open.

\medskip
\noindent\textbf{Main result.}
For any fixed polynomial-time many-one NP-complete clean-core language \(K\), proving that \(K\) has no polynomial-size family of unrestricted Boolean circuits is \emph{equivalent} to proving

\[
\NP\not\subseteq \mathsf{P/poly}.
\]

Thus a proposed ``circuit lower bound for the clean core'' is not a known weaker landing point between the current structural lemmas and \(P\ne NP\).  It is a stronger nonuniform separation.  Restricted circuit, proof-system, communication, monotone, or locality lower bounds remain useful only after an explicit model-matched transfer theorem.

\section*{Claim card}

\begin{longtable}{p{0.23\linewidth}p{0.71\linewidth}}
\toprule
\textbf{Field} & \textbf{Specification}\\
\midrule
Clean-core language & \(K=\) Monotone \(3\)-SAT-\((2,2)\): every clause has exactly three distinct variables and is all-positive or all-negative; clauses are unique; every variable occurs exactly twice positively and twice negatively.  D\"ocker proves this restriction NP-complete.  Inputs use a fixed self-delimiting indexed binary encoding; invalid strings are NO instances.\\
Encoded length & \(N=|x|\), including delimiters, signs, indices, and headers.  The characteristic function is defined on every \(x\in\{0,1\}^{N}\).\\
Circuit model & Nonuniform fan-in-two Boolean circuits over \(\{\wedge,\vee,\neg\}\); changes among finite complete bases alter size only polynomially.  Let \(C_K(N)\) be minimum circuit size for \(K\cap\{0,1\}^{N}\).\\
Exact target & 
\(
\forall c\in\mathbb N\;\forall N_0\;\exists N\ge N_0:\ C_K(N)>N^c.
\)
This is exactly \(K\notin\mathsf{P/poly}\).\\
Quantifier order & One fixed language \(K\); for every polynomial exponent and threshold, a later hard length exists.  No per-length change of machine or language is permitted.\\
Output & One decision bit.  No witness, compatibility table, separator, state transcript, or extracted object is required from the deciding circuit.\\
Load-bearing status & \textbf{OPEN}: the target is equivalent to \(\NP\not\subseteq\mathsf{P/poly}\).\\
\bottomrule
\end{longtable}

\section*{CLEANCORE-GENERAL-CIRCUIT-ENDPOINT-EQUIV1}

\begin{theorem}[\textbf{PROVED/EXACT}; independent review pending]
Let \(K\in\NP\) be NP-complete under polynomial-time many-one reductions.  Then
\[
K\notin\mathsf{P/poly}
\quad\Longleftrightarrow\quad
\NP\not\subseteq\mathsf{P/poly}.
\]
Equivalently,
\[
\left(\forall c\,\forall N_0\,\exists N\ge N_0:\ C_K(N)>N^c\right)
\quad\Longleftrightarrow\quad
\NP\not\subseteq\mathsf{P/poly}.
\]
\end{theorem}

\begin{proof}
The reverse implication is immediate because \(K\in\NP\).  For the forward closure direction, assume \(K\in\mathsf{P/poly}\), and let \(L\in\NP\).  NP-completeness gives a polynomial-time many-one reduction \(f\) satisfying
\[
x\in L\quad\Longleftrightarrow\quad f(x)\in K,
\qquad |f(x)|\le |x|^a
\]
for some constant \(a\).  A polynomial-time computation has polynomial-size circuits.  Compose a circuit computing \(f\) with the appropriate polynomial-size slice circuits for \(K\).  Variable output lengths cause no difficulty: advice for source length \(n\) may contain the target circuits for every length \(m\le n^a\), and a length selector activates the appropriate slice.  If the size-\(m\) target circuit is at most \(m^c\), their total size is \(\sum_{m\le n^a}O(m^c)=O(n^{a(c+1)})\), still polynomial.  Hence \(L\in\mathsf{P/poly}\).  Therefore \(K\in\mathsf{P/poly}\) implies \(\NP\subseteq\mathsf{P/poly}\), and contraposition completes the first equivalence.

For the second, membership in \(\mathsf{P/poly}\) means that there exist constants \(c,N_0\) with \(C_K(N)\le N^c\) for every \(N\ge N_0\); negating these quantifiers gives the displayed formula.  Finite exceptional lengths are absorbed by the circuit family.
\end{proof}

\paragraph{Exponent ledger.}
If the reduction output length is at most \(n^a\), its circuit has size at most \(n^r\), and the target slice circuit has size at most \(m^c\), the general variable-length construction above has size
\[
O\!\left(n^{\max\{r,a(c+1)\}}\right).
\]
If the reduction has one fixed output length, the sharper exponent is \(\max\{r,ac\}\).  The exponent \(c\) is not assumed known in advance; for each fixed circuit family it is a fixed constant.  Padding an encoding from \(N\) to \(M=N^p\) rescales an \(M^c\) bound to \(N^{pc}\).  This preserves polynomial versus superpolynomial size but invalidates representation-sensitive numerical exponents unless encoding is fixed.

\section*{Candidate attacked and disposition}

\paragraph{CLEANCORE-GENERAL-CIRCUIT-ENDPOINT-AS-WEAKER-BRIDGE1.}
The candidate said that a superpolynomial unrestricted-circuit lower bound for the fixed clean core might provide a weaker unconditional contradiction than the arbitrary-algorithm runtime bridge.  \textbf{FALSIFIED/RETIRED.}  The theorem above makes the proposed endpoint exactly \(\NP\not\subseteq\mathsf{P/poly}\), which implies \(P\ne NP\) and is not supplied by existing restricted-model circuit lower bounds.

\paragraph{Why assuming \(P=NP\) does not contradict known circuit theory.}
If \(P=NP\), then \(K\in P\), so its length-\(N\) functions have P-uniform polynomial-size circuits.  This is consistent with all currently known unrestricted-circuit lower bounds for explicit NP languages.  Deriving a contradiction therefore requires proving the displayed general-circuit lower bound itself or proving a valid compiler from every such circuit/decider into a restricted model with a known lower bound.  The latter compiler is exactly the missing transfer theorem; assuming it without proof hides the desired separation.

\section*{SAME\_REALIZATION and arbitrary-algorithm audit}

The equivalence theorem is language-level and nonuniform.  It does \emph{not} claim that a SAT program exposes \(H(F)\), a compatibility separator, a witness, a resolution refutation, a communication transcript, or any particular intermediate state.  It also does not infer runtime from the number of semantic states.

Any attempted implication
\[
\text{restricted representation lower bound}
\Longrightarrow
\text{runtime lower bound for every uniform SAT decider}
\]
still needs a \textbf{SAME\_REALIZATION} compiler that turns the actual decider execution into the charged restricted realization with only polynomial overhead.  This arrow is \textbf{OPEN}.  Separate external extraction, a black-box simulation, or a post hoc circuit representation is insufficient unless all queries, preprocessing, generation, random access, reconstruction, and advice are charged.

\section*{Hostile controls and finite experiment}

The bounded audit is at
\begin{center}
\path{experiments/review/clean_core_circuit_endpoint_audit.py}
\end{center}
with output
\begin{center}
\path{results/raw/review/2026-10-12_clean_core_circuit_endpoint.json}.
\end{center}
It completed in one sub-second chunk (no long computation was hidden) and checked:
\begin{itemize}
  \item all \(8^4=4096\) tuples \(a,r,c,p\in\{1,\ldots,8\}\) against the general variable-length composition rule \(\max\{r,a(c+1)\}\), the exact-length rule \(\max\{r,ac\}\), and padding rule \(pc\);
  \item 36 equality checks across direct, copy, redundant-clause, delayed, recomputed, and one-bit wrappers;
  \item 12 clause/literal-order and duplication encoding equalities and 30 padding-length checks;
  \item brute-force versus GF(2) Gaussian-elimination decisions for a Tseitin-style affine control;
  \item satisfiable and unsatisfiable 2-SAT, Horn, affine, bounded-treewidth chain, and independent-product families.
\end{itemize}

\noindent\textbf{COMPUTATIONAL only.}  These tests validate the implementation and expose wrapper/encoding sensitivity; they do not extrapolate a circuit lower bound.  The families deliberately show that \(H(F)\) may contain \(0,1,4,\) or \(8\) assignments while the observable decision is one bit and can be computed through semantically different routes.

\section*{Barrier-transfer literature preflight}

\paragraph{Obstacle 1: unrestricted-circuit endpoint.}
Razborov--Rudich identify general superpolynomial circuit lower bounds for explicit NP functions as the strong nonuniform target and give a conditional barrier for natural proofs.  Aaronson--Wigderson show that major questions including \(P\) versus \(NP\) and \(NEXP\) versus \(P/poly\) require non-algebrizing techniques.  These are not impossibility theorems for all proof methods, but they block any claim that a familiar constructive algebraic invariant automatically reaches unrestricted circuits.

\paragraph{Obstacle 2: transfer from a restricted resource to arbitrary computation.}
Fortnow--Lipton--van Melkebeek--Viglas prove genuine polynomial time--space tradeoffs for SAT on random-access machines only in explicitly bounded joint resource regimes.  Their theorem is directly transferable when its machine and resource hypotheses match; it does not yield a superpolynomial unrestricted time or circuit lower bound.  The current clean-core routes lack the required arbitrary-decider compiler.

\begin{longtable}{p{0.24\linewidth}p{0.25\linewidth}p{0.43\linewidth}}
\toprule
\textbf{Domain} & \textbf{Classification} & \textbf{Object correspondence / missing theorem}\\
\midrule
Complexity closure & directly transferable theorem & Polynomial many-one reduction circuits compose with \(K\)-slice circuits, giving the theorem above.\\
SAT time--space tradeoffs & directly transferable only under the stated RAM time and space bounds & Missing: a proof that every hypothetical polynomial-time decider lies in the charged small-space regime.\\
Monotone, bounded-depth, branching-program, proof-system lower bounds & useful structural analogy requiring a new transfer theorem & Missing: polynomial-overhead SAME\_REALIZATION compiler from every unrestricted circuit/decider to that restricted model.\\
VLSI area--time, routing congestion, communication cuts & useful structural analogy requiring a new transfer theorem & Missing: a layout/cut realization forced on arbitrary algorithms; alternate memory hierarchies and global gates are hostile controls.\\
Coding-theoretic Schur squares, quadratic relations, Hilbert functions, syzygies & useful structural analogy requiring a new transfer theorem & These may distinguish a chosen compatibility representation but do not lower-bound unrestricted circuits; hidden algebraic structure may instead enable an algorithm.\\
Control observability / reconstruction & misleading analogy without a forced late-query interface & A one-bit static decision has no obligation to preserve a state from which all coordinates can later be reconstructed.\\
\bottomrule
\end{longtable}

\section*{Post-quantum technique ledger}

\begin{longtable}{p{0.33\linewidth}p{0.58\linewidth}}
\toprule
\textbf{Technique} & \textbf{Disposition for this endpoint}\\
\midrule
Schur square; shortening/puncturing/duality & \textbf{FAILED TRANSFER} to unrestricted circuits.  Exact for a chosen code representation only.\\
Quadratic-form relation ranks; Hilbert functions; syzygies & \textbf{ANALOGY/CANDIDATE needing a new bridge}.  A distinguisher may expose exploitable structure rather than hardness.\\
Collision entropy & The Cauchy--Schwarz image bound is an \textbf{EXACT OBJECT/THEOREM TRANSFER} for an explicitly defined map, but \textbf{FAILED TRANSFER} from large image/collision entropy to arbitrary circuit size.\\
Worst-case/average-case architecture & \textbf{ANALOGY/CANDIDATE needing a new bridge}.  No unconditional average-case reduction or black-box extractor is supplied.\\
Collision extraction / special soundness & \textbf{FAILED TRANSFER} absent two conflicting openings generated from the same charged realization.\\
Hash/random-oracle methods & \textbf{FAILED TRANSFER}; conditional cryptographic binding or ideal-oracle behavior is inadmissible as an unconditional endpoint.\\
\bottomrule
\end{longtable}

\section*{Relativization, algebrization, natural-proofs, and nonuniformity audit}

The closure equivalence is definitional/reduction-theoretic and relativizes in the corresponding oracle setting when an oracle-relative complete language and reductions are fixed.  Therefore it supplies no nonrelativizing contradiction.  A general-circuit separation must confront nonuniformity directly: a different circuit is permitted at every length, but a single polynomial exponent must bound the family.  The per-length hardwired-machine quantifier inversion rejected in \texttt{POINTINF1} is not enough.  Natural Proofs applies conditionally only if the proposed property is large and constructive; algebrization applies to proof techniques, not to the elementary equivalence.  Neither barrier is being presented as an unconditional impossibility theorem.

\section*{CORPUS REVIEWED}

The run enumerated the current GitHub tree and read the current root \path{main.tex}, \path{AGENTS.md}, \path{CODEX.md}, \path{.codex/skills/pnp-research/SKILL.md}, all specifications under \path{tasks/}, current \path{main4.tex} and \path{main5.tex}, the newest verified MAIN snapshot \path{research/main/2026-10-12_0850_MAIN_SYNTHESIS-0023.tex}, and the newest verified LSPV-1 through LSPV-4 reports.  Materially relevant historical reports included \texttt{RESOURCEAXIS1}, \texttt{HEVALLB1}, \texttt{SATQUOT1}, \texttt{MAJORANTNOGO1}, \texttt{COREEQUIV1}, \texttt{POINTINF1}, \texttt{UGENCOMP1}, \texttt{LAZYACCESS1}, and \texttt{MODELBRIDGE-0023}.  The causal branch review included HEC-SAT, VSC-HEC, G-PPC-HEC, and CRS-PPC-SAT, together with cited experiment artifacts and current write-recovery governance.

Primary sources checked before substantive claims were D\"ocker's Monotone \(3\)-SAT-\((2,2)\) hardness paper, Karp's reducibility framework, Razborov--Rudich, Aaronson--Wigderson, and the SAT time--space work of Fortnow--Lipton--van Melkebeek--Viglas.

\section*{DUPLICATION / CONTRADICTION AUDIT}

\begin{itemize}
  \item \texttt{COREEQUIV1} already proves that an all-decider superpolynomial \emph{time} lower bound on a fixed NP-hard core is equivalent to \(P\ne NP\).  The present theorem is not a duplicate: it classifies the distinct nonuniform unrestricted-circuit endpoint as \(\NP\not\subseteq\mathsf{P/poly}\).
  \item \texttt{POINTINF1} rejects per-length pointwise machine minima.  The present quantifiers fix one language and require one polynomial exponent for the whole circuit family, so no hardwired machine is silently changed inside the lower-bound quantifier.
  \item Root discussions of monotone and restricted circuit lower bounds are consistent only when kept representation-specific.  Any wording suggesting that such bounds directly constrain unrestricted SAT deciders would be overstated.
  \item \path{main4.tex} gives ARC compression/counting statements conditional on unresolved all-input overlap/global consistency premises.  \path{main5.tex} gives wrapper-local exact counting with the arbitrary-solver transfer open.  This report neither redefines nor strengthens either manuscript.
\end{itemize}

\section*{WHAT CHANGED RELATIVE TO ROOT MAIN}

The root document already records that restricted circuit/proof-system lower bounds do not automatically transfer.  The new exact addition is the endpoint classification
\[
\text{superpolynomial unrestricted circuits for fixed NP-complete core}
\iff \NP\not\subseteq\mathsf{P/poly}.
\]
Accordingly, the unrestricted-circuit route is removed from the list of potentially weaker intermediate contradictions.  The highest-value surviving task is a genuinely model-matched transfer into a restricted model with a known lower bound, or the existing selected-output absolute landing theorem; both remain \textbf{OPEN}.

\section*{ADOPTED / REJECTED / RETIRED PRIOR COMPONENTS}

\begin{itemize}
  \item \textbf{ADOPTED}: \texttt{COREEQUIV1} fixed-core sufficiency; \texttt{POINTINF1} quantifier guard; \texttt{UGENCOMP1}/\texttt{LAZYACCESS1}/\texttt{MODELBRIDGE-0023} charged generation/access accounting; exact HEC no-reread results only inside their declared online model.
  \item \textbf{REJECTED}: inference from semantic compatibility count, proof length, monotone circuit size, or code invariant directly to unrestricted circuit size or arbitrary runtime.
  \item \textbf{RETIRED}: \texttt{CLEANCORE-GENERAL-CIRCUIT-ENDPOINT-AS-WEAKER-BRIDGE1}.
\end{itemize}

\section*{STALE OR OVERSTATED CLAIMS FOUND}

Any statement treating a superpolynomial \emph{general} circuit lower bound for the clean core as easier than the main separation is stale/overstated; for an NP-complete core it is exactly \(\NP\not\subseteq\mathsf{P/poly}\).  Restricted-model lower bounds may remain valid, but must be labeled by model and cannot be promoted without a transfer theorem.

\section*{Surviving target and exact missing arrow}

\[
\boxed{
\begin{gathered}
\textbf{OPEN:}\quad
\text{for every uniform deterministic polynomial-time SAT decider }A,\\
\text{construct from the same charged realization of }A
\text{ a restricted object }R_A\\
\text{with total polynomial overhead and a known model-matched lower bound,}
\end{gathered}}
\]
or else prove the selected-output absolute landing target already recorded in root main.  The construction must charge input reading, preprocessing, extraction queries, generation, access, reconstruction, and advice, retain the arbitrary solver exponent, and survive wrappers, affine elimination, alternate encodings, padding, and the one-bit output interface.  No such universal implication is proved here.

\section*{Primary references}

\begin{itemize}
  \item J. D\"ocker, ``Monotone 3-Sat-(2,2) is NP-complete,'' \url{https://arxiv.org/abs/1912.08032}.
  \item R. M. Karp, ``Reducibility Among Combinatorial Problems,'' 1972, \url{https://doi.org/10.1007/978-1-4684-2001-2_9}.
  \item A. A. Razborov and S. Rudich, ``Natural Proofs,'' JCSS 55 (1997), \url{https://doi.org/10.1006/jcss.1997.1494}.
  \item S. Aaronson and A. Wigderson, ``Algebrization: A New Barrier in Complexity Theory,'' ECCC TR08-005, \url{https://eccc.weizmann.ac.il/report/2008/005/}.
  \item L. Fortnow, R. Lipton, D. van Melkebeek, and A. Viglas, ``Time-Space Lower Bounds for Satisfiability,'' JACM 52 (2005), \url{https://pages.cs.wisc.edu/~dieter/Research/sat-ts.html}.
\end{itemize}

\bigskip
\noindent\textbf{Final status.} \textbf{PROVED/EXACT} endpoint equivalence; \textbf{FALSIFIED/RETIRED} weaker-general-circuit-bridge candidate; \textbf{COMPUTATIONAL} finite audit; \textbf{OPEN} arbitrary-decider transfer and all \(P\) versus \(NP\) consequences.  P versus NP remains open.

% END_OF_VERIFIED_REPORT::LSPV-1::CIRCEQUIV1
\end{document}
